\documentclass[10pt,a4paper]{article}
\usepackage[T1]{fontenc}
\usepackage{lmodern}
\renewcommand{\familydefault}{\sfdefault}
\usepackage[margin=1.6cm,top=1.3cm,bottom=1.3cm]{geometry}
\usepackage{microtype}
\usepackage{xcolor}
\usepackage{eurosym}
\usepackage{enumitem}
\usepackage{tabularx}
\usepackage{booktabs}
\usepackage[most]{tcolorbox}
\usepackage[colorlinks=true,urlcolor=blue!50!black]{hyperref}
\pagestyle{empty}
\setlength{\parindent}{0pt}
\setlength{\parskip}{0.35em}
\setlist{itemsep=0.1em,topsep=0.2em,leftmargin=1.2em}

\definecolor{cblue}{HTML}{2A78D6}
\definecolor{corange}{HTML}{EB6834}
\definecolor{cgray}{HTML}{52514E}

\newcommand{\head}[1]{\par\smallskip{\large\bfseries\color{cblue}#1}\par\vspace{0.15em}}
\newtcolorbox{stat}{colback=cblue!5,colframe=cblue!60,boxrule=0.6pt,arc=2pt,left=4pt,right=4pt,top=3pt,bottom=3pt,
  before upper=\raggedright}

\begin{document}

{\LARGE\bfseries AI in customer service: triage now, auto-replies later}\par
\vspace{0.2em}
{\color{cgray}Solution brief for a retail bank \textbullet{} Albin Oruglica \textbullet{} September 2026 \textbullet{}
Based on a working prototype and 13{,}083 real banking questions}

\head{The situation}
A bank's service team answers thousands of short, often incomplete messages: \emph{``my card hasn't arrived''},
\emph{``why was I charged twice?''}. Answering them needs information from several systems. An AI assistant could
handle many of them instantly, but a confident wrong answer to a customer is costly, and some cases (suspected
fraud, disputes, anti-money-laundering checks) must never be decided by a machine. The business question is
\textbf{which messages AI should handle, and at what point it should bring in a human.}

\head{What was built and tested}
A prototype that reads each message, identifies the request, looks up the customer's (simulated) card,
account and transaction data, and then answers, asks one clarifying question, or hands over to an agent with a
case summary. It was tested on a public set of real banking questions, with fictional customer data. Two AI
components were compared: a small, cheap \textbf{classifier} that sorts messages, and a \textbf{large language
model} (Claude) that looks up data and writes replies.

\head{Key findings}
\begin{tcbraster}[raster columns=3,raster equal height,raster column skip=6pt]
\begin{stat}\textbf{\large Sorting works.}\\ The classifier identified the customer's request correctly in
\textbf{89\%} of cases, at practically zero cost, and knows when it is unsure.\end{stat}
\begin{stat}\textbf{\large Safety is a setting.}\\ At a strict setting, \textbf{none of 440} fraud, dispute and
money-laundering cases slipped past to automation; a looser setting automates more but lets about 1\% through.\end{stat}
\begin{stat}\textbf{\large Replies are not ready.}\\ Even with the request correctly identified, the AI fully
resolved only \textbf{1 in 5} messages, and about \textbf{1 in 5} of its answers contained an error.\end{stat}
\end{tcbraster}

\begin{itemize}
  \item \textbf{Where the value is.} Sorting messages and routing them to the right place could be worth up to about
  \textbf{\euro2--2.5 per message} against today's fully manual handling (assumed \euro5 per message). This is a ceiling:
  it assumes every sorted message then gets a correct reply.
  \item \textbf{Why not let the AI reply on its own.} When the AI's own replies were checked end to end, the fully
  automated system cost about the same as, or slightly more than, human-only handling (\euro5.1--6.0 vs.\ \euro5.00
  per message). The errors were plausible-sounding: a real rule applied to the wrong product, an invented policy.
  Simple automatic checks do not catch them.
  \item \textbf{The right threshold depends on the bank's own costs} (what a wrong answer and an agent's time really
  cost, and how much risk is acceptable), and it must be re-measured whenever the model changes. It cannot be set
  once from a formula.
\end{itemize}

\head{Recommendation}
\begin{tabularx}{\linewidth}{@{}>{\bfseries\color{cblue}}p{2.6cm}X@{}}
Phase 1: now & \textbf{AI triage and drafting, humans send.} The classifier sorts every message and routes
fraud, dispute and anti-money-laundering cases straight to specialists. For the rest, the AI prepares a draft
reply and a case summary, and the agent reviews, edits and sends. Low risk; no AI text reaches a customer
unchecked.\\[0.3em]
Phase 2: later & \textbf{Auto-send for a narrow set of simple requests}, only once measured reply accuracy on the
bank's own data beats the break-even (roughly fewer than 1 wrong answer in 6 at these costs), backed by a proper
knowledge base of products and procedures, and with the strict safety setting.\\
\end{tabularx}

\textbf{Controls from day one:} fixed, auditable rules for cases AI must never decide; monitoring of how many
sensitive cases the sorting lets through; the AI sees only the customer data it needs; a data-processing
agreement and legal basis for any external AI service (GDPR), or an in-house model if data must stay internal.

\head{Proposed next step: a 6-week pilot with one service team}
Run Phase 1 on live traffic and measure what the prototype could not: \textbf{agent handling time with and
without AI drafts}, the share of drafts sent unchanged, sorting accuracy on the bank's real mix of requests, and
the number of sensitive cases misrouted. These four numbers decide whether Phase 1 pays for itself and when Phase 2
is justified.

\vfill
{\footnotesize\color{cgray} Evidence: prototype, code and all results in the project repository (see its README;
every number can be reproduced with one command). Tested on the public BANKING77 dataset (PolyAI, CC BY 4.0) with invented customer
data and an 11-line fictional product fact sheet; results are indicative, not a forecast for a specific bank. Cost
figures use assumptions (\euro5 per manual message, \euro30 per wrong answer, \euro300 per misrouted sensitive
case) and are tested for sensitivity in the repository.}

\end{document}
